\section{MCU and digital interfaces}
\label{sec:esp32-mcu}

An ESP32-S3-WROOM-1 module supplies the control logic. It drives the four BCD
decoder/drivers (\cref{sec:bcd-driver}). It reads the flyback power-good signal
(\cref{sec:hv-house-pgood}). It also supplies a serial port for firmware and
console access. The module is the main 3.3~V load.
\Cref{sec:3v3-converter} uses its supply current.

\begin{table}[htbp]
\centering
\caption{MCU and interface component designators}\label{tab:esp32-designators}
\begin{tabular}{|l|l|}
\hline
\textbf{Designator} & \textbf{Function} \\
\hline
U501 & ESP32-S3-WROOM-1 module \\
\hline
C501, C502 & Module local decoupling \\
\hline
C503 & Reset (EN) node capacitor \\
\hline
J501 & Serial header \\
\hline
U502 & Inbound serial buffer (header to MCU) \\
\hline
U503 & Outbound serial buffer (MCU to header) \\
\hline
C504, C505 & Buffer supply decoupling \\
\hline
R502, R504 & Buffer input pull-ups \\
\hline
R503 & Outbound series resistor \\
\hline
\end{tabular}
\end{table}

\subsection{ESP32-S3-WROOM-1 module}
\label{sec:esp32}

Source: \texttt{esp32-s3-wroom-1\_wroom-1u\_datasheet\_en.pdf}, Espressif. The
datasheet covers variants with different flash and PSRAM configurations.
\textbf{Use a variant without PSRAM.} Octal-interface variants use GPIO35 to
GPIO37 for the data bus. This design needs these pins as general-purpose
outputs. Flash capacity does not affect the design.

{\footnotesize
\begin{longtable}{|L{4.3cm}|c|c|c|c|L{3.0cm}|c|}
\caption{ESP32-S3-WROOM-1 electrical characteristics}\label{tab:esp32}\\
\hline
\textbf{Parameter} & \textbf{Min} & \textbf{Typ} & \textbf{Max} & \textbf{Units} & \textbf{Conditions} & \textbf{Page} \\
\hline
\endfirsthead
\caption[]{ESP32-S3-WROOM-1 electrical characteristics (continued)}\\
\hline
\textbf{Parameter} & \textbf{Min} & \textbf{Typ} & \textbf{Max} & \textbf{Units} & \textbf{Conditions} & \textbf{Page} \\
\hline
\endhead
VDD33 recommended operating range & \num{3.0} & \num{3.3} & \num{3.6} & \si{\volt} & & 27 \\
\hline
VDD33 absolute maximum & \num{-0.3} & & \num{3.6} & \si{\volt} & & 27 \\
\hline
External supply current capability required & \num{0.5} & & & \si{\ampere} & ``Current delivered by external power supply'' & 27 \\
\hline
Wi-Fi TX peak (highest of all published conditions) & & & \num{355} & \si{\milli\ampere} & 802.11b, 1~Mbps, @20.5~dBm; rated at 100\% TX duty cycle & 28 \\
\hline
Wi-Fi TX, 802.11n HT40 MCS7 @17~dBm & & & \num{285} & \si{\milli\ampere} & & 28 \\
\hline
Wi-Fi RX, HT20 & & \num{95} & & \si{\milli\ampere} & Peripherals disabled, CPU idle & 28 \\
\hline
Wi-Fi RX, HT40 & & \num{97} & & \si{\milli\ampere} & Peripherals disabled, CPU idle & 28 \\
\hline
Bluetooth LE TX peak & & & \num{344} & \si{\milli\ampere} & @20.0~dBm (max power) & 28 \\
\hline
Bluetooth LE RX & & \num{93} & & \si{\milli\ampere} & & 28 \\
\hline
Modem-sleep, lowest & & \num{13.2} & & \si{\milli\ampere} & WAITI, dual-core idle, 40~MHz & 29 \\
\hline
Modem-sleep, lowest, peripheral clocks enabled & & \num{18.8} & & \si{\milli\ampere} & WAITI, dual-core idle, 40~MHz & 29 \\
\hline
Modem-sleep, highest & & \num{91.7} & & \si{\milli\ampere} & Dual-core, 128-bit access, 240~MHz & 29 \\
\hline
Modem-sleep, highest, peripheral clocks enabled & & \num{107.9} & & \si{\milli\ampere} & Dual-core, 128-bit access, 240~MHz & 29 \\
\hline
Light-sleep & & \num{240} & & \si{\micro\ampere} & Wi-Fi and VDD\_SPI powered down, GPIOs Hi-Z & 30 \\
\hline
Deep-sleep (various ULP configurations) & \num{7} & & \num{190} & \si{\micro\ampere} & & 30 \\
\hline
Power off & & \num{1} & & \si{\micro\ampere} & EN held low & 30 \\
\hline
EN recommended external network & \multicolumn{4}{L{5.0cm}|}{\(R=10\)~k\(\Omega\), \(C=1\)~\si{\micro\farad}} & & 41 \\
\hline
EN reset-release threshold & \(0.75\times V_{DD}\) & & & & & 27 to 28 \\
\hline
EN reset threshold & & & \(0.25\times V_{DD}\) & & & 27 to 28 \\
\hline
EN/strapping setup time & \num{0} & & & \si{\milli\second} & & 14 \\
\hline
EN/strapping hold time & \num{3} & & & \si{\milli\second} & & 14 \\
\hline
GPIO0 (boot strap) default state & \multicolumn{4}{L{5.0cm}|}{Weak pull-up; high = SPI boot (normal), low = download boot} & & 13 to 14, 27 \\
\hline
ESD rating, HBM & & & \num{2000} & \si{\volt} & $\pm$, either polarity & 47 \\
\hline
ESD rating, CDM & & & \num{500} & \si{\volt} & $\pm$, either polarity & 47 \\
\hline
Flash current adder (modem-sleep, 80~Mbit/s, SPI 2-line) & & \num{10} & & \si{\milli\ampere} & & 29 \\
\hline
GPIO \(V_{IH}\) & \(0.75\times V_{DD}\) & & & & \SI{25}{\celsius} & 27 \\
\hline
GPIO \(V_{IL}\) & & & \(0.25\times V_{DD}\) & & \SI{25}{\celsius} & 27 \\
\hline
GPIO \(V_{OH}\) & \(0.8\times V_{DD}\) & & & & High-impedance load & 27 \\
\hline
GPIO \(V_{OL}\) & & & \(0.1\times V_{DD}\) & & High-impedance load & 27 \\
\hline
GPIO source current & & \num{40} & & \si{\milli\ampere} & \(V_{DD}=3.3\)~V, PAD\_DRIVER=3 & 27 \\
\hline
GPIO sink current & & \num{28} & & \si{\milli\ampere} & \(V_{DD}=3.3\)~V, PAD\_DRIVER=3 & 27 \\
\hline
GPIO internal weak pull-up/down (each) & & \num{45} & & \si{\kilo\ohm} & & 27 \\
\hline
\end{longtable}
}

\paragraph{Data not specified.} The module datasheet does not give the VDD33
brownout or power-on-reset threshold. It does not give a GPIO absolute-maximum
voltage other than the recommended \(V_{DD}+0.3\)~V limit. It also gives no
aggregate GPIO-current limit. It does not characterize \(V_{OL}\) or \(V_{OH}\)
at a specified load current. The chip-level datasheet contains some of these
values. \Cref{sec:esp32-bcd} states the inferred output-stage resistance.

\paragraph{GPIO overvoltage.} The recommended input maximum is
\(V_{DD}+0.3\)~V. A voltage above approximately \SI{3.6}{\volt} can forward-bias
the input clamp (inferred). This limit applies to each external interface.
\Cref{sec:esp32-uart} supplies the necessary isolation.

\paragraph{Supply-current requirement.} Espressif requires an external supply
that can supply at least \SI{0.5}{\ampere}. The maximum specified Wi-Fi transmit
current is \SI{355}{\milli\ampere}. A supply sized only for idle current can
cause an undervoltage during transmission. \Cref{sec:3v3-converter} uses the
\SI{0.5}{\ampere} requirement.

\subsection{Logic buffer}
\label{sec:lvc1g17}

The serial interface needs a single-gate noninverting buffer. It must have an
overvoltage-tolerant Schmitt input and a push-pull output. It must also specify
partial-power-down operation. Source: \texttt{74LVC1G17.pdf}, Nexperia Rev.
16.1, 2024-09. The table uses the \(V_{CC}=3.0\)~V data across
\SIrange{-40}{85}{\celsius}.

{\footnotesize
\begin{longtable}{|L{4.4cm}|c|c|c|c|L{3.4cm}|c|}
\caption{74LVC1G17-class buffer characteristics}\label{tab:lvc1g17}\\
\hline
\textbf{Parameter} & \textbf{Min} & \textbf{Typ} & \textbf{Max} & \textbf{Units} & \textbf{Conditions} & \textbf{Page} \\
\hline
\endfirsthead
\caption[]{74LVC1G17-class buffer characteristics (continued)}\\
\hline
\textbf{Parameter} & \textbf{Min} & \textbf{Typ} & \textbf{Max} & \textbf{Units} & \textbf{Conditions} & \textbf{Page} \\
\hline
\endhead
Supply voltage \(V_{CC}\), recommended & \num{1.65} & & \num{5.5} & \si{\volt} & & 5 \\
\hline
Input voltage \(V_I\), recommended & \num{0} & & \num{5.5} & \si{\volt} & & 5 \\
\hline
Output voltage \(V_O\), recommended, active & \num{0} & & \(V_{CC}\) & \si{\volt} & Active mode & 5 \\
\hline
Output voltage \(V_O\), recommended, powered down & \num{0} & & \num{5.5} & \si{\volt} & Power-down mode; \(V_{CC}=0\)~V & 5 \\
\hline
Supply voltage \(V_{CC}\), absolute maximum & \num{-0.5} & & \num{6.5} & \si{\volt} & & 4 \\
\hline
Input voltage \(V_I\), absolute maximum & \num{-0.5} & & \num{6.5} & \si{\volt} & & 4 \\
\hline
Output voltage \(V_O\), absolute maximum, active & \num{-0.5} & & \(V_{CC}+0.5\) & \si{\volt} & Active mode & 4 \\
\hline
Output voltage \(V_O\), absolute maximum, powered down & \num{-0.5} & & \num{6.5} & \si{\volt} & Power-down mode; \(V_{CC}=0\)~V & 4 \\
\hline
Output current \(I_O\), either direction & & & \num{50} & \si{\milli\ampere} & \(V_O=0\)~V to \(V_{CC}\) & 4 \\
\hline
Output clamping current \(I_{OK}\), either direction & & & \num{50} & \si{\milli\ampere} & \(V_O>V_{CC}\) or \(V_O<0\)~V & 4 \\
\hline
Positive-going threshold \(V_{T+}\) & \num{1.29} & \num{1.5} & \num{1.71} & \si{\volt} & \(V_{CC}=3.0\)~V & 6 \\
\hline
Negative-going threshold \(V_{T-}\) & \num{0.88} & \num{1.0} & \num{1.24} & \si{\volt} & \(V_{CC}=3.0\)~V & 6 \\
\hline
Hysteresis \(V_H\) & \num{0.31} & \num{0.5} & \num{0.64} & \si{\volt} & \(V_{CC}=3.0\)~V & 6 \\
\hline
\(V_{OH}\), high-level output voltage & \(V_{CC}-0.1\) & & & \si{\volt} & \(I_O=-100\)~\si{\micro\ampere}; \(V_{CC}=1.65\)~V to \num{5.5}~V & 5 \\
\hline
\(V_{OH}\), high-level output voltage & \num{2.3} & & & \si{\volt} & \(I_O=-24\)~mA; \(V_{CC}=3.0\)~V & 5 \\
\hline
\(V_{OL}\), low-level output voltage & & & \num{0.1} & \si{\volt} & \(I_O=100\)~\si{\micro\ampere}; \(V_{CC}=1.65\)~V to \num{5.5}~V & 5 \\
\hline
\(V_{OL}\), low-level output voltage & & & \num{0.55} & \si{\volt} & \(I_O=24\)~mA; \(V_{CC}=3.0\)~V & 5 \\
\hline
\(I_I\), input leakage current, either polarity & & & \num{1} & \si{\micro\ampere} & \(V_I=5.5\)~V or GND; \(V_{CC}=0\)~V to \num{5.5}~V & 5 \\
\hline
\(I_{OFF}\), power-off leakage current, either polarity & & & \num{2} & \si{\micro\ampere} & \(V_I\) or \(V_O=5.5\)~V; \(V_{CC}=0\)~V & 5 \\
\hline
\(I_{CC}\), supply current & & & \num{4} & \si{\micro\ampere} & \(V_I=5.5\)~V or GND; \(I_O=0\)~A & 5 \\
\hline
\(C_I\), input capacitance & & \num{5} & & \si{\pico\farad} & & 5 \\
\hline
\(t_{pd}\), propagation delay & \num{0.7} & \num{3.0} & \num{5.5} & \si{\nano\second} & \(V_{CC}=3.0\)~V to \num{3.6}~V & 8 \\
\hline
Latch-up performance & \num{250} & & & \si{\milli\ampere} & & 1 \\
\hline
ESD rating, HBM & \num{2000} & & & \si{\volt} & JS-001 class 2 & 1 \\
\hline
ESD rating, CDM & \num{1000} & & & \si{\volt} & JS-002 class C3 & 1 \\
\hline
\end{longtable}
}

A substitute buffer must meet four requirements. Its input must accept
\SI{5.5}{\volt} independently of \(V_{CC}\). It must specify power-off leakage
with \SI{5.5}{\volt} on each terminal. It must have a Schmitt input. It must
have a push-pull output.

\paragraph{Input and output tolerance.} The input accepts \SI{5.5}{\volt} in
all supply states. During operation, the output accepts only voltages to
\(V_{CC}\). It accepts \SI{5.5}{\volt} only when \(V_{CC}=0\). A driven output
above \(V_{CC}\) can conduct into the supply (inferred). Therefore, the outbound
line has a series resistor. The inbound line does not need one.

\subsection{BCD drive interface}
\label{sec:esp32-bcd}

Sixteen GPIOs drive the BCD inputs, with four GPIOs for each driver. The MCU
supply spans \SIrange{3.0}{3.6}{\volt}. The drivers use 5~V TTL input levels.
Thus, both logic levels require verification. The checks use the
\SI{3.0}{\volt} minimum MCU supply.

\paragraph{High level.} The MCU must hold a driver input above the TTL
\(V_{IH}\) of \SI{2.0}{\volt} while sourcing that input's high-level current,
at most \SI{80}{\micro\ampere} on the B, C, and D inputs:
\begin{align*}
  V_{OH} &\geq 0.8\,V_{DD} = 0.8\times\SI{3.0}{\volt} = \SI{2.40}{\volt}
  \quad\text{against}\quad V_{IH} = \SI{2.0}{\volt},
\end{align*}
The margin is \SI{0.40}{\volt} at minimum supply. It increases to
\SI{0.64}{\volt} at \SI{3.3}{\volt}. The \(V_{OH}\) specification uses a
high-impedance load. The driver draws at most \SI{80}{\micro\ampere}, which is
$0.2\,\%$ of the GPIO source-current rating. Thus, the interface meets the
fixed TTL high-level threshold.

\paragraph{Low level.} The MCU must hold a driver input below the TTL
\(V_{IL}\) of \SI{0.8}{\volt} while sinking that input's low-level current, at
most \SI{3.2}{\milli\ampere} on the B, C, and D inputs. The module's
\(V_{OL}\) specification is given only into a high-impedance load, and
\SI{3.2}{\milli\ampere} is not that, so the requirement is restated as a bound
on the pin's effective pull-down impedance:
\begin{align*}
  R_{\mathrm{pd}} &\leq \frac{\SI{0.8}{\volt}}{\SI{3.2}{\milli\ampere}}
    = \SI{250}{\ohm}.
\end{align*}
The \SI{28}{\milli\ampere} sink rating implies an effective pull-down
resistance below this limit (inferred). If the pin gives \(0.1\,V_{DD}\) at the
rated current, the resistance is approximately \SI{11}{\ohm}. It then gives
approximately \SI{35}{\milli\volt} at \SI{3.2}{\milli\ampere}. The interface
accepts any pull-down resistance to \SI{250}{\ohm}.

\paragraph{Aggregate low-level current.} Worst case, all sixteen lines low:
\begin{align*}
  I_{\mathrm{BCD,sink}} &= 4 \times \left(\SI{1.6}{\milli\ampere}
    + 3\times\SI{3.2}{\milli\ampere}\right) = \SI{44.8}{\milli\ampere},
\end{align*}
sourced from the 5~V rail through the drivers' input structures and returned
through the MCU's ground, so it places no load on the 3.3~V rail. The drivers'
own supply-current figure already includes it, since its test condition is
inputs grounded (\cref{sec:bcd-driver}), the same all-inputs-low state. The 5~V
budget of \cref{sec:5v-load} therefore counts it exactly once. Per pin,
\SI{3.2}{\milli\ampere} against a \SI{28}{\milli\ampere}
rating is unremarkable; the module datasheet publishes no aggregate GPIO
current limit against which to check the \SI{44.8}{\milli\ampere} total, so
none is claimed.

\paragraph{Cathode mapping.} The decoder outputs do not connect to cathodes in
numeric order. A numeric connection order would cause excessive trace crossings
on a two-layer board. Therefore, each cathode connects to the most direct driver
output. Firmware uses a mapping table for each tube. Read each mapping from the
schematic. The four tubes can have different mappings.

\paragraph{State before firmware operation.} MCU GPIOs are high impedance after
reset. Thus, the driver inputs float until firmware configures them. A floating
TTL input usually reads high (inferred). Binary 1111 turns all decoder outputs
off. Noise can change this state. However, all reachable decoder states are
electrically safe. Firmware then drives the inputs to defined levels.

\paragraph{GPIO count.} Twenty-five GPIOs are required: sixteen BCD outputs,
two serial signals, one power-good input, one display-inhibit output, one USB
switch output, three HPI signals, and one port-controller fault input. The
native USB D+ and D$-$ pins are additional dedicated peripheral pins.
Assignment must avoid holding any
boot-strapping pin at a level that selects an unintended boot mode through
the \SI{3}{\milli\second} strapping hold time after reset release. The count is
far below the module's available pins, so the constraint is on which pins are
chosen.

\subsection{Serial header}
\label{sec:esp32-uart}

J501 exposes the MCU's UART (universal asynchronous receiver/transmitter) for
firmware download and console access, driven by a commodity USB-to-serial
adapter. The header's supply pin feeds the 3.3~V converter's input and is
treated in \cref{sec:3v3-diodeor}; this subsection covers the two signal
lines.

\paragraph{Adapter signal voltage.} USB-to-serial adapters can use 3.3~V or 5~V
signals. A 5~V signal exceeds the MCU input limit by approximately
\SI{1.4}{\volt}. A series resistor limits clamp current but does not limit pin
voltage. Therefore, no external adapter signal can connect directly to an MCU
pin.

\paragraph{Buffers.} Both lines pass through 74LVC1G17-class buffers
(\cref{sec:lvc1g17}): U502 inbound from the header to the MCU, U503 outbound
from the MCU to the header. Both use the 3.3~V supply.
\begin{itemize}[nosep]
  \item \textbf{Inbound.} The input accepts up to \SI{5.5}{\volt} for any
    permitted supply voltage. The output is limited to $V_{CC}$
    (\cref{tab:lvc1g17}). A \SI{3.3}{\volt} supply converts 5~V adapter signals
    to MCU-compatible levels.
  \item \textbf{Outbound.} U503 isolates the MCU's transmit pin from the
    connector. A wiring fault reaches the buffer output.
  \item \textbf{Unpowered board.} The buffers' power-off leakage is
    \SI{2}{\micro\ampere} with \SI{5.5}{\volt} on either terminal at
    \(V_{CC}=0\). This limits current into the unpowered 3.3~V rail.
  \item \textbf{Signal integrity.} The Schmitt-trigger input tolerates the
    cable edges. Hysteresis is at least \SI{0.31}{\volt} at a
    \SI{3.0}{\volt} supply. Maximum propagation delay is
    \SI{5.5}{\nano\second}.
\end{itemize}
Each buffer has a \SI{2}{\kilo\volt} human-body-model rating. Use a local
\SI{100}{\nano\farad} supply capacitor for each buffer.

\paragraph{Idle-state pull-ups.} Pull both buffer inputs up to the 3.3~V rail.
R502 controls U502. R504 controls U503. This produces the high serial idle
state. Each pull-up must hold its input above the
\SI{1.71}{\volt} worst-case positive-going threshold against the input's
\SI{1}{\micro\ampere} leakage, bounding it at
\((3.0-1.71)/\SI{1}{\micro\ampere}\approx\SI{1.3}{\mega\ohm}\). Set
$R502=R504=\SI{10}{\kilo\ohm}$. Each resistor draws
\SI{330}{\micro\ampere} when its line is low.

\paragraph{Outbound series resistor.} U503's output connects to the header.
Limit a ground-short current to \SI{25}{\milli\ampere}:
\begin{align*}
  R_{503} &\geq \frac{\SI{3.3}{\volt}}{\SI{25}{\milli\ampere}}
    = \SI{132}{\ohm}.
\end{align*}
An external \SI{5}{\volt} drive produces \SI{8}{\milli\ampere} of clamp
current through \SI{150}{\ohm}. The clamp-current rating is
\SI{50}{\milli\ampere}. Set $R503=\SI{150}{\ohm}$.

\paragraph{Adapter input level.} U503 supplies at least approximately
\SI{3.2}{\volt} at light load. The adapter must have a TTL-compatible input
threshold. A 5~V CMOS input can require \SI{3.5}{\volt} and is not compatible.

\paragraph{Adapter requirement.} The interface accepts 3.3~V and 5~V adapter
signals. The adapter input must have TTL-compatible thresholds. Its supply pin
must also meet \cref{sec:3v3-diodeor}.

\subsection{Native USB}
\label{sec:esp32-usb}

The ESP32-S3 USB Serial/JTAG peripheral connects to the USB-C receptacle through
the switch in \cref{sec:pd-data}. It supports firmware download, console, and
debug functions. The MCU can start from a \SI{5}{\volt} host port.

The link uses USB full speed. Both USB blocks support full speed only
(\ds{esp32-s3-wroom-1\_wroom-1u\_datasheet\_en.pdf}{21,22}). Thus, the data rate
is \SI{12}{\mega\bit\per\second}. \Cref{sec:pd-data} uses this limit when it
connects the controller and MCU to one pair. The circuit cannot support USB high
speed because the controller pin capacitance exceeds that interface budget.

The serial header remains available when native USB is disabled. Native USB
requires an operating input stage and USB peripheral.

\paragraph{USB buffers.} Do not add logic buffers to D+ or D$-$. USB~2.0 uses
a bidirectional pair at the MCU logic level. \Cref{sec:pd-esd} specifies the
connector ESD protection.

\paragraph{Power-up pulses.} During start-up, both pins have two high pulses.
Each pulse is approximately \SI{60}{\micro\second}. They occur in
\SI{3.2}{\milli\second} and \SI{2}{\milli\second} windows. The switch in
\cref{sec:pd-data} stays open during this interval. It isolates the pulses from
the receptacle.

\paragraph{Switch control.} One GPIO drives the data switch's enable. Two
requirements bound which pin it may be:
\begin{itemize}[nosep]
  \item It must not be a strapping pin. The strapping pins are sampled at reset
    (\cref{sec:esp32-reset}), and a pin that both selects boot mode and holds
    the data switch open would couple the two.
  \item Its reset state must be high impedance, so that \des{R}{USBSW}
    (\cref{sec:pd-data}) holds the switch open through reset and through any
    interval in which firmware is not driving it. The controller owns the pair
    at attach and must keep owning it until the MCU deliberately takes it.
\end{itemize}
Assign the switch control and display inhibit to general-purpose pins. Their
external resistors set safe states while the pins are high impedance.

\subsection{Power-good input}
\label{sec:esp32-pgood}

The flyback PGOOD output is open-drain. \des{R}{HVPG} pulls it up to the 3.3~V
rail (\cref{sec:hv-house-pgood}). An MCU GPIO reads the node. PGOOD releases
when the flyback output, supply, and UVLO conditions are valid
(\ds{lm5156h.pdf}{25}).

The high level is approximately \(V_{DD}\). The input threshold is
\(0.75\,V_{DD}\). The controller sinks \SI{33}{\micro\ampere} for a low level.
If the MCU pull-up is enabled, the sink current rises to
\SI{106}{\micro\ampere}. The controller sink limit is
\SI{1}{\milli\ampere}.

The MCU and pull-up use the same supply. Thus, the signal cannot drive an
unpowered MCU. The pull-up draws \SI{33}{\micro\ampere} during start-up and
fault conditions.

\subsection{Reset and boot configuration}
\label{sec:esp32-reset}

EN must stay below \(0.25\,V_{DD}\) while the supply rises
(\cref{sec:esp32}). The Espressif reference circuit uses an RC delay. This
design uses the 3.3~V converter power-good output instead
(\cref{sec:3v3-house}). \des{R}{3V3PG}, a \SI{100}{\kilo\ohm} resistor,
pulls up the node.

EN connects to the converter output through \des{R}{3V3PG}. Therefore, EN
cannot rise faster than the supply. PG holds EN low until the output exceeds
\SI{94.5}{\percent} of its programmed value. It then releases EN. No second
pull-up is necessary. PG sinks \SI{33}{\micro\ampere} while it is asserted.

If the rail falls below \SIrange{90}{94.5}{\percent}, PG asserts again. It
holds the module in reset until regulation returns. An RC network cannot give
this continuous supervision.

\paragraph{Reset capacitor.} \(C503 = \SI{100}{\nano\farad}\) filters contact
bounce from the mechanical reset switch. With \des{R}{3V3PG}, it gives a
\SI{10}{\milli\second} time constant. It also filters noise on the EN node. A
\SI{1}{\micro\farad} capacitor would slow manual reset assertion without a
necessary benefit.

Download-mode entry requires GPIO0 held low across a reset release and for the
\SI{3}{\milli\second} strapping hold time that follows
(\cref{sec:esp32}). The serial header carries no modem-control lines, so no
automatic entry path exists and the design requires a manual means of
asserting GPIO0 and EN together. GPIO0's internal weak pull-up selects normal
boot with no external component, so only the asserting path is a requirement,
not a pull-up.

\subsection{Summary}\label{sec:esp32-summary}

The analysis uses three inferred values. The MCU GPIO pull-down resistance is
derived from its \SI{28}{\milli\ampere} sink rating. Serial connector and cable
capacitance is at most \SI{50}{\pico\farad}. The selected adapter has
TTL-compatible input thresholds.

{\footnotesize
\begin{longtable}{|p{0.23\textwidth}|p{0.335\textwidth}|p{0.335\textwidth}|}
\caption{MCU and digital interfaces: specification summary}\label{tab:esp32-summary}\\
\hline
\textbf{Parameter} & \textbf{Derived requirement} & \textbf{Specification / result} \\
\hline
\endfirsthead
\caption[]{MCU and digital interfaces: specification summary (continued)}\\
\hline
\textbf{Parameter} & \textbf{Derived requirement} & \textbf{Specification / result} \\
\hline
\endhead
\hline
\endlastfoot
BCD high level & $\geq\SI{2.0}{\volt}$ into
  $\leq\SI{80}{\micro\ampere}$ (TTL $V_{IH}$) & \SI{2.40}{\volt} at the
  \SI{3.0}{\volt} supply floor \\
\hline
BCD low level & $\leq\SI{0.8}{\volt}$ while sinking
  \SI{3.2}{\milli\ampere}; pin impedance $\leq\SI{250}{\ohm}$ &
  $\approx\SI{11}{\ohm}$ implied by the \SI{28}{\milli\ampere} sink rating \\
\hline
BCD aggregate sink & Not applicable & \SI{44.8}{\milli\ampere} from the 5~V rail,
  already inside the drivers' $I_{CC}$ \\
\hline
GPIO count & 16 BCD $+$ 2 serial $+$ 1 power-good $+$ 1 display inhibit $+$
  1 USB switch $+$ 3 HPI $+$ 1 fault & 25 control, serial, and BCD GPIOs;
  dedicated native USB pins are additional \\
\hline
Serial buffers U502/U503 & input recommended to \SI{5.5}{\volt} at any
  $V_{CC}$; power-off leakage specified at $V_{CC}=0$; Schmitt input;
  push-pull output & 74LVC1G17 class, 3.3~V supply \\
\hline
Buffer decoupling C504/C505 & \SI{100}{\nano\farad} per logic package &
  \SI{100}{\nano\farad} $\times 2$ \\
\hline
Input pull-ups R502/R504 & hold input above \SI{1.71}{\volt} against
  \SI{1}{\micro\ampere} leakage: $\leq\SI{1.3}{\mega\ohm}$ &
  \SI{10}{\kilo\ohm}; \SI{330}{\micro\ampere} when driven low \\
\hline
Outbound series R503 & bound short-circuit current below the
  \SI{50}{\milli\ampere} output rating: $\geq\SI{132}{\ohm}$ &
  \SI{150}{\ohm}; \SI{22}{\milli\ampere} shorted, \SI{8}{\milli\ampere}
  clamping at a 5~V external drive \\
\hline
Adapter signaling level & unconstrained & 3.3~V and 5~V both in
  specification \\
\hline
Adapter input thresholds & TTL-referenced ($\SI{2.0}{\volt}$ $V_{IH}$) &
  driven at $\approx\SI{3.2}{\volt}$ \\
\hline
Power-good input & levels within $0.25/0.75\,V_{DD}$ & reads
  $\approx V_{DD}$ high, near ground low; \SI{33}{\micro\ampere} load \\
\hline
Reset (EN) node & EN low during ramp, released after &
  3.3~V converter PG drives it, terminated by \des{R}{3V3PG};
  released at \SI{94.5}{\percent} of the rail \\
\hline
C503 & swallow reset-switch bounce; noise immunity on a
  \SI{100}{\kilo\ohm} node & \SI{100}{\nano\farad};
  $\tau=\SI{10}{\milli\second}$ against \des{R}{3V3PG} \\
\hline
Download-mode entry & GPIO0 low across reset $+$ \SI{3}{\milli\second} &
  manual assertion required; no modem-control lines on the header \\
\hline
\end{longtable}
}
